\section{Theory and its limits}\label{sec:theory}
The local proximal identities must be separated from convergence of a full
splitting algorithm and from finite-precision numerical acceptance. Each
admitted scalar or directional proximal problem has a proper closed convex
objective plus a positive quadratic metric. It therefore has a unique
minimizer. The original QP itself may still be infeasible, unbounded or
nonunique. A local row formula cannot remove those global possibilities.
Standard proximal and monotone-operator results provide the appropriate
framework for the distinction~\cite{prox,rockafellarppa,lions,eckstein}.

For the grouped split define a graph-constrained function
\begin{equation}\label{eq:paper-graph-function}
 F(x,z)=\tfrac12x^\top P x+q^\top x+
            \iota_{\{z=Ax\}}(x,z),
\end{equation}
and let $G$ contain the hard intervals and the admitted soft costs.
The fixed block metric is $\mathcal M=\diag(S,H)$ with $S\succ0$ and
positive diagonal $H$. The graph subproblem is uniquely defined because
its reduced matrix $P+S+A^\top H A$ is positive definite. This is
numerical regularization of a proximal subproblem, not an undisclosed
addition to the scientific objective. The target problem remains the
original constrained convex optimization model.

Write $J_F,J_G$ for the proximal resolvents in this fixed metric and
$R_F=2J_F-I$, $R_G=2J_G-I$ for their reflectors. The native consensus
state uses the scaled dual $d$ and the following four updates:
\begin{align}\label{eq:paper-consensus-updates}
 u^{k+1}&=J_F(v^k-d^k),\\
 h^{k+1}&=\alpha u^{k+1}+(1-\alpha)v^k,\nonumber\\
 v^{k+1}&=J_G(h^{k+1}+d^k),\nonumber\\
 d^{k+1}&=d^k+h^{k+1}-v^{k+1}.\nonumber
\end{align}
Define $s^k=v^k+d^k$ after a complete step, so $v^k=J_G(s^k)$.
Substitution with this specific state gives
\begin{equation}\label{eq:paper-relaxed-dr}
 s^{k+1}=\left(1-\frac\alpha2\right)s^k+
                    \frac\alpha2 R_F R_G(s^k),\qquad0<\alpha<2.
\end{equation}
Under an attained saddle point and the stated closed-convex assumptions,
the standard fixed-metric splitting theory supplies convergence of the
appropriate shadow variables. The relaxation range makes the displayed
operator averaged in the relevant metric. Its fixed points are connected
to primal/dual solutions through the resolvents, rather than being
identified with raw numerical coordinates without that recovery.
This interpretation follows established splitting results~\cite{admm,eckstein}.

In original row variables, writing $S$ for the positive primal metric
and $y=Hd_z$, the same graph solve and relaxed updates take the form
\begin{align}\label{eq:paper-native-updates}
 (P+S+A^\top HA)\widetilde x^{k+1}
   &=Sx^k-q+A^\top(Hz^k-y^k),\\
 x^{k+1}&=\alpha\widetilde x^{k+1}+(1-\alpha)x^k,\nonumber\\
 z^{k+1}&=\prox_g^H\left(\alpha A\widetilde x^{k+1}
                       +(1-\alpha)z^k+H^{-1}y^k\right),\nonumber\\
 y^{k+1}&=y^k+H\left(\alpha A\widetilde x^{k+1}
                       +(1-\alpha)z^k-z^{k+1}\right).\nonumber
\end{align}
The price function gains no relaxation factor. The grouped native
implementation fixes a finite relaxation in the stated open interval
and rejects adaptive penalty updates, automatic OSQP scaling, polishing,
directional sensitivities and directional code generation. Its explicit
row normalization is separate: $\bar A=DA$, $\bar\ell=D\ell$,
$\bar u=Du$, $\bar c=c/D$, $y=D\bar y$ and
$H_{\rm physical}=D\bar H D$. Setup chooses
$D_i=\max(c_{L,i},c_{U,i})$ on soft rows and one on hard rows.
This identifies the physical metric; it does not claim that a
previously unnormalized iteration has the same trajectory.

It is not a rate theorem for every admitted semidefinite QP. Stronger
contraction or linear-rate conclusions need additional quantitative
assumptions~\cite{giselsson}. Adaptive changes in the row metric or the
relaxation parameter produce a different iteration and require their
own conditions. The empirical original protocol's documented adaptation
is evaluated numerically; the fixed-metric argument is not presented
as a proof of every adaptive schedule in every implementation. Likewise,
iteration-limited candidate vectors remain subject to the external
acceptance rule regardless of their distance from a recent iterate.

The diagonal assumption in the translated proximal identity cannot be
discarded. Consider $C=(-\infty,0]\times\R$, $\Psi=\iota_{\{0\}}$,
$v=(1,0)$ and
\begin{equation}\label{eq:paper-dense-counterexample}
 H=\begin{bmatrix}2&1\\1&2\end{bmatrix}.
\end{equation}
The residual formula would return the Euclidean box projection $(0,0)$.
However, the true $H$-metric projection is $(0,1/2)$: fixing the active
first coordinate at zero, differentiation of the quadratic in the
second coordinate gives $-1+2z_2=0$. The metric is positive definite,
and the residual indicator satisfies coordinate contraction monotonicity,
yet the identity fails. Positive definiteness alone is therefore not
an adequate replacement for a positive diagonal metric in that theorem.

Equivalent formulations provide additional mathematical controls. A
directional upper row is $c_{U,i}A_ix\leq c_{U,i}u_i$, while a lower
row is $-c_{L,i}A_ix\leq-c_{L,i}\ell_i$. Their positive departures
produce the same directional peak. A compact epigraph introduces
$t_g\geq0$ and bounds every scaled departure by $t_g$, charging
$a_gt_g+\tfrac12b_gt_g^2$. Nonnegative prices permit a tight epigraph
at an optimum, with harmless nonuniqueness when the price is free.
These are exact objective representations, not weakened reference
problems. They motivate independent numerical comparisons without
providing a theorem about which representation is fastest.

Finally, original-space acceptance includes assumptions not proved by
an algebraic identity alone. Sparse source conversion, finite precision,
dual recovery and physical units can be wrong even when the intended
formula is correct. Independent lifting, local Fenchel relations,
hard-row complementarity and reciprocal scaling controls test those
implementation boundaries. Such checks supply finite evidence on actual
inputs. They do not replace the mathematical domain restrictions, nor
establish operational safety, global priority or a complete-operation
timing gain.
